\documentclass[UTF8]{ctexart}
\usepackage{ctex}

\usepackage{mydoc}
\usepackage{xjtlumathstyle}
\usepackage{xjtlumaththm}

\usepackage{listings}
\usepackage{xcolor}

\usepackage{graphicx}
\usepackage{doc}

\lstset{language=TeX,
basicstyle=\small,
backgroundcolor=\color{black!20},
breaklines=true,
breakindent=0pt% disable indention for line breaks
}

\usepackage{hyperref}
\hypersetup{colorlinks=true}

\newcommand{\gitrepo}{\url{https://github.com/Little-He-Guan/A-XJTLU-Math-Club-LaTeX-Template}}
\title{宏包 xjtlumath 文档}
\author{贺官羽铭\\  {\footnotesize\gitrepo}}
\date{\today}

\newcommand{\mycprt}{版权所有 \copyright{} 贺官羽铭 2021-\the\year.}
\begin{document}
\pagenumbering{gobble}
\maketitle

\vfill
\section*{版权}
\mycprt{} 此文档被The MIT License保护.

\clearpage
\pagenumbering{roman}
\tableofcontents

\part{基本信息}
\pagestyle{headings}
\section{这宏包是干嘛的？}
我最初这个 \LaTeX{} 包最初是为西浦数学社的资料部使用的。 然而我发现我的大学缺乏学生的 \LaTeX{} 模板，因此我决定将其扩展为其它学生也可以使用的模板。

该宏包包含一些用于数学文档的有用命令和环境，并重新定义了一些现有样式，使其适合我们大学的样式。 

\section{历史}
2020年，我加入了西浦数学社的资料部。该部门制作有关数学的资料并将其分发给学生以帮助他们。在我加入时人们正在用Microsoft Word 和MathType 准备资料。我看到这是一个用\LaTeX 的力量来增强资料的好机会，所以我建议使用\LaTeX 来准备文档。

建议通过了，我们的团队开始转移我们的工作环境。在这个过程中，我写了一个小包，即是这个包的前身，在我们的资料中使用。我们部门很多人对\LaTeX{}几乎一无所知，我承诺新的准备过程不会很困难，我设法定义了整体程序，分享了一些关于\LaTeX 的必要知识，并教他们如何使用此宏包。

我承认我当时有点不负责任。我确实描述了工作程序是如何以及我的包应该如何使用，但我较少分享关于 \LaTeX 的内容。我的包和工作流简化了一些概念，因此团队不需要关心 \LaTeX{} 的某些部分，例如documentclass 和preamble。然而 \LaTeX{} 本身与 Word 仍有很大不同。我低估了我的同事学习基本的 \LaTeX{} 技巧的难度，所以这对他们中的一些人来说是一种折磨。后来我有些意识到这个问题，但我被自己的业务所占据，并没有给予足够的重视。

最终的结果基本是是成功的，而且大多已经足够好了。但我为我的同事感到难过，因为这件事对他们来说实际上并不那么容易。另外，我最近参加了西浦数学俱乐部15周年纪念活动，这在某种程度上增强了我对数学俱乐部的归属感。出于这些原因，我决定将来为西浦的学生留下一些东西。我重写了整个包，并添加了一系列教程，涵盖了\LaTeX{} 的基础知识以及制作材料，以便资料部门的人在 \LaTeX 中制作资料的工作量会大大减少。教程采用故事的形式，因为我认为这种方法为学习\LaTeX{} 增添了更多乐趣。

除此之外，西浦的学生甚至其他 \LaTeX{} 用户可能会发现这个包和文档有用。所以它的最终形式被决定为一个带有文档的模板，在我的个人帐户下作为 GitHub repository 发布。这是因为我们的数学社目前没有 GitHub 帐户。一旦它建立了一个账户，我将把这个存储库转移到它的官方帐户。

\section{中文版历史}
由于XJTLU的 University Communication Language是英语，所以此文档最开始被用英语写作。此中文版来源于对英文版的翻译。

%\part{Tutorials} inside Tutorial 1
\part{Tutorials}
\pagenumbering{arabic}
\section{Tutorial 1: Ashley's First Material}
Ashley joined the material department recently. Now he is assigned to write a part of the material that covers basic calculus. He is very excited because this will be his first material as well as his first try on \LaTeX. First, he needs to know how to use \LaTeX{} and the package on his PC. In this tutorial, we will follow Ashley and see what he learns in the material preparation.

\subsection{Installation and Configuration}
To use \LaTeX{}, Ashley needs a \TeX{} distribution (and probably an editor) installed on his computer. There are several popular distributions listed on the \LaTeX{} official website: \url{https://www.latex-project.org/get/}. The installation and configuration of these distributions are quite easy and Ashley completed them in a few minutes.

Then, Ashley wants to have an editor for writing \LaTeX{} documents. He learns that \TeX studio is a good one, so he downloads and installs it.

What Ashley needs to do now is to obtain a copy of this package and install it. (TBD)

\subsection{Basic Writings}
Ashley, waiting breathlessly to begin his first work, clicks open the Templates folder and navigates into the material-book folder. He opens encapsulation.tex and notices that it is like this:
\begin{lstlisting}
... Some code ...

\input{chapters.tex}

... Some code ...
\end{lstlisting}

As the file name suggests, this is for making a material that -is to be published as a book. A book contains some chapters, and in the material writing, each chapter is a specific part of the material topic. For example, Ashley's work is part of the calculus material, so he needs to start a chapter for his work. He opens chapter.tex and finds a blank document. He then adds the following command to start his first chapter and runs \LaTeX{} on encapsulation.tex to check the output (On the left is the result of his command(s), appearing in a different font family, and on the right is his commands, with a gray background):

\begin{miniexammar}{.45\textandmarginlen}{\fakesectioningdef{1}{Key points in calculus}}
\begin{lstlisting}
\chapter{Key points in calculus}
\end{lstlisting}
\end{miniexammar}
\verb=\chapter{...}= is a \LaTeX{} \emph{command}, which starts with a \verb=\=%
. The words being wrapped with the curly brackets \verb={}= form an \emph{argument} for the command. At this place, they are the caption of the chapter. Ashley notices that \LaTeX{} automatically enlarges the font and makes them bold. Being different with many other typesetting software, \LaTeX{} only needs the logical idea of what to do (e.g. there will be a chapter named xxx at some place), and it will control the appearance for the author.

Ashley is satisfied with that result. He then types some of paragraphs. In \LaTeX{}, paragraphs are separated by one blank line.

\begin{miniexammar}{.5\textandmarginlen}{
\fakesectioningdef{1}{Key points in calculus}%
What does Ashley write in these paragraphs? Well, in fact, I don't know. You may find him and ask him yourself.
		
\hspace{1.5em}But, wait a minute, how do I find Ashley when he doesn't really exist? Well, this is a good question.}
\begin{lstlisting}
\chapter{Key points in calculus}
What does Ashley write in these paragraphs? Well, in fact, I don't know. You may find him and ask him yourself.
		
But, wait a minute, how do I find Ashley when he doesn't really exist? Well, this is a good question.
\end{lstlisting}
\end{miniexammar}
Ashley notices that paragraphs are automatically indented by \LaTeX{}. Yet he also notices that the paragraph directly under the chapter is not indented.

Apart from chapter, \LaTeX{} also provides these sectioning commands:
\begin{itemize}
\item section
\item subsection
\item subsubsection
\item paragraph
\item subparagraph
\end{itemize}
You may notice that paragraph is included here. In fact, the sectioning command \verb=\paragraph= generates a title for the paragraphs like other sectioning commands.

When Ashley turns to the whole output, he sees that his chapter appears in the table of contents:

\begin{miniexammar}{.7\textandmarginlen}{
\faketoc
\fakecontentsline{1}{\textbf{Key points in calculus}}{1}
}
\begin{lstlisting}
\chapter{Key points in calculus}
\end{lstlisting}
\end{miniexammar}
\LaTeX{} generates the table of contents for all sectioning commands\footnote{whose depth are below the table of contents (toc) depth. The actual process of how toc is generated is a bit complex, so I will not cover it here}. For this reason, Ashley has to run \LaTeX{} twice on encapsulation.tex for a correct table of contents.

When Ashley once writes a very long caption for a section, the table of contents become quite ugly. To solve this, Ashley can specify a short form of the section that is used in the table of contents, as shown below:
\begin{miniexammar}{.7\textandmarginlen}{
\faketoc
\fakecontentsline{1}{\textbf{A Short Name}}{1}
}
\begin{lstlisting}
\chapter[A Short Name]{A very very long Caption}
\end{lstlisting}
\end{miniexammar}

\subsection{Adjust fonts} \label{subsec:fonts}
So far, Ashley knows how to instruct \LaTeX{} to do some basic things. Although he feels full of energy and is writing at full speed, he soon encounters some problems. Ashley wants to emphasize some keywords such as ``limit''. He later learns that the command \verb=\emph= instructs \LaTeX{} to emphasize the text passed to it, as shown below.

\begin{parexammar}{.38\textandmarginlen}{
Calculus is the study of \emph{limits}.
}
\begin{lstlisting}
Calculus is the study of \emph{limits}.
\end{lstlisting}
\end{parexammar}

This thing is quite \LaTeX, as Ashley only tells \LaTeX{} to emphasize it, and has no control of how \LaTeX{} does it. Although most of times it is enough, Ashley wants more. He wonder how can one explicitly control the appearance of the texts, since \LaTeX{} cannot cover all needs in every situations. \LaTeX{} does provide certain default operations on fonts, and Ashley can use them to control the size, family, and style of his texts.

\begin{parexammar}{.45\textandmarginlen}{
Ashley can tell \LaTeX{} to adjust the font size like:
{\tiny very very small} {\scriptsize the size of scripts} {\footnotesize the size of foot notes} {\small small font} {\normalsize just being normal} {\large a bit bigger} {\Large large text} {\LARGE very big} {\huge huge} {\Huge damn huge}
}
\begin{lstlisting}
Ashley can tell \LaTeX{} to adjust the font size like:
{\tiny very very small} {\scriptsize the size of scripts} {\footnotesize the size of foot notes} {\small small font} {\normalsize just being normal} {\large a bit bigger} {\Large large text} {\LARGE very big} {\huge huge} {\HUGE damn huge}
\end{lstlisting}
\end{parexammar}
This time Ashley sees something different from the command \verb=\emph{}=. The texts here are inside curly brackets, and the commands are just given inside the brackets along with the texts. Something enclosed by a pair of curly brackets is said to be inside a \emph{group}. Commands called inside the group influence the whole group.

Apart from the sizes, Ashley is also able to control the style and family of fonts. As shown in Table \ref{tab:stdfontcmds}, \LaTeX{} provides these commands to control font style and family:
\begin{table}[htbp]
\begin{tabular}{lll}
Command & Used in a group & Action\\
\hline
\verb=\textrm{...}= & \verb={\rmfamily...}= & {Text in \textrm{Roman} family} \\
\verb=\textsf{...}= & \verb={\sffamily...}= & {Text in \textsf{sans serif} family} \\
\verb=\texttt{...}= & \verb={\ttfamily...}= & {Text in \texttt{typewriter} family} \vspace{.15cm}\\
\verb=\textmd{...}= & \verb={\mdseries...}= & {Text in \textmd{medium} series} \\
\verb=\textbf{...}= & \verb={\bfseries...}= & {Text in \textbf{bold} series} \vspace{.15cm}\\
\verb=\textup{...}= & \verb={\upshape...}= & {Text in \textup{upright} shape} \\
\verb=\textit{...}= & \verb={\itshape...}= & {Text in \textit{italic} shape} \\
\verb=\textsl{...}= & \verb={\slshape...}= & {Text in \textsl{slanted} shape} \\
\verb=\textsc{...}= & \verb={\scshape...}= & {Text in \textsc{small caps} shape} \vspace{.15cm}\\
\verb=\emph{...}= & \verb={\em...}= & {Text \emph{emphasized}}\vspace{.15cm}\\
\verb=\textnormal{...}= & \verb={\normalfont...}= & {Text in default font}
\end{tabular}
\caption{Standard font-changing commands and declarations}
\label{tab:stdfontcmds}
\end{table}

Many of these commands provide both an used-in-group version and a normal version. Ashley can choose which version to use depend on his needs.

Ashley also wants to learn how to change the color of fonts. He is surprised that this package documentation does not provide such a description. After contacting with the package author, he learns that this topic is not covered because the materials are printed in black and white, so it would be nearly useless to change colors.

\subsection{Typing Mathematical Formulae}
It comes the most exciting part of Ashley's work --- typing formulae! Even though he had little experience in \LaTeX{} before, he already learned that \LaTeX{} produces high-quality math formulae, as he previously saw at some sites like Math Stack Exchange (\url{https://math.stackexchange.com/}) and ZhiHu (\url{https://www.zhihu.com/}).

Because of his previous experience on these sites, he knows a little bit about how to write formulae in \LaTeX{}.

Generally, formulae in \LaTeX{} are classified into two types: \emph{inline} and \emph{displayed}. Formulae of the former type are enclosed in a pair of dollar sign: \verb=$...$=, while formulae of the latter type are enclosed in a pair of double dollar sign: \verb=$$...$$=. These delimiters are the original \TeX{} ones. \LaTeX{} provides additionally two pairs of delimiters for inline and displayed math, respectively: \verb=\(...\)= and \verb=\[...\]=. In fact, the \TeX{} shorthand \verb=$$...$$= for displayed math should be avoided, as it may lead to strange problems in \LaTeX{}.

As their name suggest, an inline formula is in a line of texts, while a displayed formula is displayed outside of the main texts.

\begin{miniexammar}{.55\textandmarginlen}%
{
The derivative of a function $f$ can be written as $f'$, or as
\[
\frdt{x}
\]
}
\begin{lstlisting}
The derivative of a function $f$ can be written as $f'$, or as
\[
\frdt{x}
\]
\end{lstlisting}
\end{miniexammar}
Here, Ashley uses a command provided by xjtlumath: \verb=\frdt=. This command takes two argument, where the first one is optional with the default value $f$. They represent the function and the variable, respectively. When Ashley wants to typeset the derivative of $g$ to $x$, he will write 
\verb=\frdt[g]{x}=, which gives $\frdt[g]{x}$. The optional argument is passed to the command in the embrace of a pair of two square brackets. As careful as Ashley, you may notice that in inline mode, a formula is shrunk to some extent so that it can fit in one line.

Ashley is quite satisfied with the rendering effect. He then writes some equations and texts. But when he wants to reference a equation, a problem arises. How does he address the equation he wants to reference? Here Ashley is introduced with another way of giving a displayed math: the equation environment.

\begin{miniexammar}{.5\textandmarginlen}%
{
The fundamental theorem of calculus can be expressed in the form of Equation \ref{eq:fundthmcal}
\begin{equation}
\int_a^b f(x) \dx = F(b) - F(a) \label{eq:fundthmcal}
\end{equation}
}
\begin{lstlisting}
The fundamental theorem of calculus can be expressed in the form of Equation \ref{eq:fundthmcal}
\begin{equation}
\int_a^b f(x) \dx = F(b) - F(a) \label{eq:fundthmcal}
\end{equation}
\end{lstlisting}
\end{miniexammar}

Ashley understands the equation by his previous knowledge: the underscore (\verb=_=) introduces the subscript ($a$) to the integral symbol $\int$ (\verb=\int=), while the caret (\verb=^=) introduces the superscript ($b$) to it. The command \verb=\dx= is provided by xjtlumath. If Ashley types dx directly, the result will be like this $\int f(x) dx$, which is ugly. \verb=\dx= refines the result. Note that \verb=\dx= can only be used to represent the integral variable, because it adds a little space before it. If Ashley needs to use other variables, he needs to use the command \verb=\dd=. xjtlumath also provides the same facility for multiple integrals.

\begin{miniexammar}{.6\textandmarginlen}%
{
\[
\int f(t) \dd t,\quad
\iint f(x,y) \dxdy,\quad
\iiint f(x,y,z) \dxdydz
\]
\[
\iint f \drdt,\quad
\iiint f \dzdrdt,\quad
\iiint f \drdtdp
\]
}
\begin{lstlisting}
\[
\int f(t) \dd t,\quad
\iint f(x,y) \dxdy,\quad
\iiint f(x,y,z) \dxdydz
\]
\[
\iint f \drdt,\quad
\iiint f \dzdrdt,\quad
\iiint f \drdtdp
\]
\end{lstlisting}
\end{miniexammar}

However, Ashley knows nothing about cross-referencing in \LaTeX{} as well as the equation environment. Let's explain them to Ashley. In \LaTeX{}, an environment is started by the \verb=\begin= command and ends at the \verb=\end= command. The name of the environment is passed to the pair of commands. The equation environment gives a displayed math equation that is \emph{counted}. The \verb=\label= command at the end of the equation catches the counter as well as some other information like its location and store it in the label represented by the name given to \verb=\label=. To use the label, Ashley needs the \verb=\ref= command, which prints the counter\footnote{and in addition generates a clickable hyperlink, which on click navigates to the location of the equation, and which is the effect of the hyperref package loaded by the template.}.

The equation environment gives a numbered counted, while \verb=\[\]= doesn't. Using a label inside this causes the label to be directed to another counter, so use a label only when that thing is counted.

\subsection{Space Management}
The space management in \LaTeX{} is a bit more complex than just typing white spaces. Ashley is a careful person, he soon finds that the space after a sentence is a little bit larger than the space between words (you may zoom in the .pdf file to see this). \LaTeX{} decides a space as a space at the end of a sentence if
\begin{enumerate}
\item A full stop (.) or right quotation mark (') is immediately followed by the space, and
\item if it is a full stop being followed, the letter immediately before the full stop is in lowercase.
\end{enumerate}
Most sentences end according to the above rules, though there are some exceptions. For example, \LaTeX{} may take a Mr.~as the sign of a sentence ending, and thus produces wrong spacing. Under such circumstances, Ashley needs to configure \LaTeX{} manually by \verb=~= and \verb=\@=.

\begin{parexammar}{.5\textandmarginlen}%
{
In another world, Mr.~Ashley was once loved by Miss Scarlett. In this world, Mr.~Ashley has a PC\@. He loves programming on his PC.
}
\begin{lstlisting}
In another world, Mr.~Ashley was once loved by Miss Scarlett. In this world, Mr.~Ashley has a PC\@. He loves programming on his PC.
\end{lstlisting}
\end{parexammar}
Ashley is happy as he learned how to manage the spaces. Yet soon he finds another problem.

\begin{parexammar}{.5\textandmarginlen}%
{
Some commands like \LaTeX seems to eat the space after it.
}
\begin{lstlisting}
Some commands like \LaTeX seems to eat the space after it.
\end{lstlisting}
\end{parexammar}
To solve this problem, Ashley needs to add an empty group (\verb={}=) at the end of the command.

As for mathematical formulae, things become different. \LaTeX{} ignores all white spaces in math mode, whether it is inline or displayed. To add extra spaces, Ashley has to use the commands shown in Table \ref{tab:spacemath}.
\begin{table}[htbp]
\begin{center}
\begin{tabular}{ll}
Command & Effect (approximately) \\
\hline
\verb=\,= & $\frac{3}{18}$ quad (\showwidth{.166666em}) \vspace{5pt}\\
\verb=\:= & $\frac{4}{18}$ quad (\showwidth{.222222em}) \vspace{5pt}\\
\verb=\;= & $\frac{5}{18}$ quad (\showwidth{.277777em}) \vspace{5pt}\\
\verb=\ = (\verb=\= followed by a space) & a space \vspace{5pt}\\
\verb=\quad= & Width of `M' in current font (\showwidth{1em}) \vspace{5pt}\\
\verb=\qquad= & 2 quad (\showwidth{2em})
\end{tabular}
\end{center}
\caption{Spacing in Math Mode}
\label{tab:spacemath}
\end{table}

\subsection{Lists and Other Environments}
Now Ashley has learned about dealing with texts, he continues his writing. Soon he has to break again as he is working on a list. At first, he hard-codes the list like this:

\begin{parexammar}{.4\textandmarginlen}{
1. Something

2. Something

3. Something
}
\begin{lstlisting}
1. Something

2. Something

3. Something
\end{lstlisting}
\end{parexammar}
Yet this solution looks rather silly. In addition, if Ashley wants to change the number, he needs to do it manually. He wonders if \LaTeX{} has some more convenient way to do it.

Fortunately there is. \LaTeX{} provides several environment to deal with lists. One example is shown below.

\begin{miniexammar}{.4\textandmarginlen}{
\begin{enumerate}
\item Something
\item Something
\item Something
\end{enumerate}
}
\begin{lstlisting}
\begin{enumerate}
\item Something
\item Something
\item Something
\end{enumerate}
\end{lstlisting}
\end{miniexammar}

Now Ashley has nearly everything he needs to know. He is content with what he has written and feels happy. The only thing that matters for him is that he wants to show theorems, definitions, and other things in a more fancy fashion so that his readers can focus on these.

xjtlumath provides some fancy environments just for this purpose.

\begin{miniexammar}{.6\textandmarginlen}{
\begin{definition}[Absolute Convergence]
An infinite series $\infseries{n}{0}$ is said to be absolutely convergent iff
\[
\sum_{n=0}^\infty |a_n|
\]
converges
\end{definition}
}
\begin{lstlisting}
\begin{definition}[Absolute Convergence]
An infinite series $\infseries{n}{0}$ is said to be absolutely convergent iff
\[
\sum_{n=0}^\infty |a_n|
\]
converges
\end{definition}
\end{lstlisting}
\end{miniexammar}

xjtlumath loads amsthm, which gives the proof environment. When Ashley uses this environment to write proofs, he finds that a q.e.d.~sign appears at the end of the environment.
\begin{miniexammar}{.5\textandmarginlen}{
Now we prove the mean value theorem for definite integrals. That is, for a continuous function $f$ that is bounded on $[a,b]$, the definite integral $\int_a^b f(x) \dx = f(c)(b-a)$, where $c \in [a,b]$. 
\begin{proof}
Let $m,M$ be the infimum and supermum of $f([a,b])$, respectively.
Therefore, $m\le f \le M$, and
\[
\int_a^b m \dx \le \int_a^b f(x) \dx \le \int_a^b M \dx
\]
, which gives
\begin{align}
m(b-a) \le &\int_a^b f(x) \dx \le M(b-a) \nonumber\\
m \le &\frac{\int_a^b f(x) \dx}{b-a} \le M \label{eq:meanvalint}
\end{align}

Since that $f$ is continuous, it can reach every value between the infimum and supermum of its range. That is, $\exists c \in [a,b], f(c) = \int_a^b f(x) \dx$. Substitute $f(c)$ back to equation \ref{eq:meanvalint} gives what the theorem states.
\end{proof}
}
\begin{lstlisting}[basicstyle=\footnotesize]
Now we prove the mean value theorem for definite integrals. That is, for a continuous function $f$ that is bounded on $[a,b]$, the definite integral $\int_a^b f(x) \dx = f(c)(b-a)$, where $c \in [a,b]$. 
\begin{proof}
Let $m,M$ be the infimum and supermum of $f([a,b])$, respectively.
Therefore, $m\le f \le M$, and
\[
\int_a^b m \dx \le \int_a^b f(x) \dx \le \int_a^b M \dx
\]
, which gives
\begin{align}
m(b-a) \le &\int_a^b f(x) \dx \le M(b-a) \nonumber\\
m \le &\frac{\int_a^b f(x) \dx}{b-a} \le M \label{eq:meanvalint}
\end{align}

Since that $f$ is continuous, it can reach every value between the infimum and supermum of its range. That is, $\exists c \in [a,b], f(c) = \int_a^b f(x) \dx$. Substitute $f(c)$ back to equation \ref{eq:meanvalint} gives what the theorem states.
\end{proof}
\end{lstlisting}
\end{miniexammar}

Ashley is able to control where the q.e.d.~sign appears by using the \verb=\qedhere= command provided by proof. If this command is given before, then it will not appear at the end.
\begin{miniexammar}{.5\textandmarginlen}{
\begin{proof}
Some words...
\[
a+b=c \qedhere
\]
Som words...
\end{proof}
}
\begin{lstlisting}
\begin{proof}
Some words...
\[
a+b=c \qedhere
\]
Som words...
\end{proof}
\end{lstlisting}
\end{miniexammar}

Besides of definition and proof, Ashley is also able to use theorem, proposition, corollary, lemma, axiom, and example. Except for proof, these environments own their individual counters, and Ashley can simply use label to reference them.

However, one colleague of Ashley, Chao, complains that the English caption of these environments conflicts with the Chinese material he is writing. For this purpose, xjtlumath additionally provides Chinese version environments in replace of them in Chinese materials. Appending a `c' to the environments' names gives the names of the Chinese version environments.

These facilities greatly help Ashley in his material preparation, and he will finish his work soon...

\section{教程2: Delilah和复杂的数学公式}
Delilah 正在研究有关线性代数的材料的一部分。 随着工作的进行，她将获得处理 \LaTeX{} 中复杂数学公式的能力，尤其是 xjtlumath 中加载的 ams 包提供的那些方法。 

\subsection{多行公式}
线性方程组是线性代数的基本部分。 当 Delilah 尝试输入一组方程时，她遇到了一个问题。 在预定义的 \verb=\[\]= 和环境equation中，她找不到开始新行的选项。 甚至 \LaTeX{} 的换行选项，如 \verb=\\= 和 \verb=\newline= 在那里也不起作用。 一个方程组当然不应该放在一行，那她现在该怎么办？ 后来她了解到环境 \emph{aligned} 旨在允许方程组在多行中对齐：
\begin{miniexammar}{.4\textandmarginlen}{
\[
\begin{aligned}
x+y &= 1\\
x-y &= 2
\end{aligned}
\]
}
\begin{lstlisting}
\[
\begin{aligned}
x+y &= 1\\
x-y &= 2
\end{aligned}
\]
\end{lstlisting}
\end{miniexammar}

这里，与符号 \verb=&= 用在方程要对齐的符号之前。 换行符 \verb=\\= 开始一行新的等式。 请注意，此处不能使用其他换行操作。

Delilah 喜欢这个结果，但她觉得方程式太孤单了。 她认为为他们添加一个大花括号会安慰他们。 \LaTeX{} 支持将分隔符放在一组事物之前和之后的语法。
\begin{miniexammar}{.4\textandmarginlen}{
\[
\left\{
\begin{aligned}
x+y &= 1\\
x-y &= 2
\end{aligned}
\right.
\]
}
\begin{lstlisting}
\[
\left\{
\begin{aligned}
x+y &= 1\\
x-y &= 2
\end{aligned}
\right.
\]
\end{lstlisting}
\end{miniexammar}
\verb=\left= 命令定义要放在左边的东西，而\verb=\right= 命令定义要放在右边的东西。 Delilah 不想把任何东西放在右边，所以写了 \verb=.=表示空。

对于一组不需要对齐的方程，或者对于一个太长而不能放在一行中的单个方程，没有对齐的环境 \emph{gathered} 是更好的选择：
\begin{miniexammar}{.45\textandmarginlen}{
\[
\begin{gathered}
\cos {z} = 1 - \frac{z^2}{2!} + \frac{z^4}{4!} - \frac{z^6}{6!} + \cdots \\
= \sum_{n=0}^\infty {\frac{(-1)^n z^{2n}}{(2n!)}}
\end{gathered}
\]
}
\begin{lstlisting}
\[
\begin{gathered}
\cos {z} = 1 - \frac{z^2}{2!} + \frac{z^4}{4!} - \cdots \\
= \sum_{n=0}^\infty {\frac{(-1)^n z^{2n}}{(2n!)}}
\end{gathered}
\]
\end{lstlisting}
\end{miniexammar}

写了几组方程后，Delilah 想引用其中的一组。 她使用equation环境而不是 \verb=\[\]=，但发现方程是作为一个整体编号的。
\begin{miniexammar}{.4\textandmarginlen}{
\begin{equation}
\begin{aligned}
x+y &= 1\\
x-y &= 2
\end{aligned}
\end{equation}
}
\begin{lstlisting}
\begin{equation}
\begin{aligned}
x+y &= 1\\
x-y &= 2
\end{aligned}
\end{equation}
\end{lstlisting}
\end{miniexammar}
所以她很难在一个组中引用单个方程。 amsmath 为此提供环境 align：

\begin{miniexammar}{.4\textandmarginlen}{
\begin{align}
x+y &= 1\\
x-y &= 2
\end{align}
}
\begin{lstlisting}
\begin{align}
x+y &= 1\\
x-y &= 2
\end{align}
\end{lstlisting}
\end{miniexammar}

如果她不想为单个方程编号，则需要在该行的末尾附加 \verb=\nonumber=。
\begin{miniexammar}{.4\textandmarginlen}{
\begin{align}
x+y &= 1\\
z&=10 \nonumber\\
x-y &= 2
\end{align}
}
\begin{lstlisting}
\begin{align}
x+y &= 1\\
z &= 10 \nonumber\\
x-y &= 2
\end{align}
\end{lstlisting}
\end{miniexammar}

如果没有 'ed' 后缀，\emph{gather} 也是一个独立的环境，可以完成gather的工作。 但是普通版本和“ed”版本之间有一个主要区别。 Delilah 发现不可能在对齐或聚集之前再次放置括号，因为它们不需要但被数学环境包围。 此外，它们的宽度固定为文本的宽度，而它们的“ed”版本可以是任何宽度。如果没有 'ed' 后缀，\emph{gather} 也是一个独立的环境，可以完成聚集的工作。 但是普通版本和“ed”版本之间有一个主要区别。 Delilah 发现不可能在对齐或聚集之前再次放置括号，因为它们不需要但被数学环境包围。 此外，它们的宽度固定为文本的宽度，而它们的“ed”版本可以是任何宽度。

与 equation 环境一样，它们的星号版本默认不给出数字。
\begin{miniexammar}{.4\textandmarginlen}{
\begin{align*}
x+y &= 1\\
x-y &= 2
\end{align*}
}
\begin{lstlisting}
\begin{align*}
x+y &= 1\\
x-y &= 2
\end{align*}
\end{lstlisting}
\end{miniexammar}

Delilah 能够将多组方程放在一个列中，只需在组之间添加\verb=&=号即可。
\begin{miniexammar}{.4\textandmarginlen}{
\begin{align*}
x+y &= 1  & a+b &= 3\\
x-y &= 2  & a-b &= 4
\end{align*}
}
\begin{lstlisting}
\begin{align*}
x+y &= 1  & a+b &= 3\\
x-y &= 2  & a-b &= 4
\end{align*}
\end{lstlisting}
\end{miniexammar}
align 自动调整方程组之间的空间。

\subsection{矩阵}
矩阵对于线性代数至关重要，因为它们表示从一个向量空间到另一个特定基底的线性映射。 此外，系数矩阵和增广矩阵便于操作线性方程。

amsmath 提供了多种输入矩阵的环境。
\begin{miniexammar}{.5\textandmarginlen}{
\[
\begin{bmatrix}
1&2&3&4\\
5&6&7&8\\
9&10&11&12\\
13&14&15&16
\end{bmatrix}
\]
}
\begin{lstlisting}
\[
\begin{bmatrix}
1&2&3&4\\
5&6&7&8\\
9&10&11&12\\
13&14&15&16
\end{bmatrix}
\]
\end{lstlisting}
\end{miniexammar}
环境 pmatrix、Bmatrix、vmatrix 和 Vmatrix 分别产生分隔符 \verb=()=、\verb={}=、\verb=||= 和 \verb=|| ||=。

为了在inline模式下使用矩阵，Delilah 使用环境 smallmatrix，它在 amsmath 中没有 p,b,B,v,V 版本，因为决定分隔符是作者的责任。
\begin{miniexammar}{.4\textandmarginlen}{
The matrix $\left(\begin{smallmatrix} a&b\\c&d \end{smallmatrix}\right)$ is so small and cute!
}
\begin{lstlisting}
The matrix $\left(\begin{smallmatrix} a&b\\c&d \end{smallmatrix}\right)$ is so small and cute!
\end{lstlisting}
\end{miniexammar}

当 Delilah 尝试将分数放入矩阵中时，她发现有些烦人事
。\begin{miniexammar}{.3\textandmarginlen}{
\[
\begin{bmatrix}
1&\frac{1}{2}&\frac{1}{3}\\
1&\frac{1}{4}&\frac{1}{5}\\
\end{bmatrix}
\]
}
\begin{lstlisting}
\[
\begin{bmatrix}
1&\frac{1}{2}&\frac{1}{3}\\
1&\frac{1}{4}&\frac{1}{5}\\
\end{bmatrix}
\]
\end{lstlisting}
\end{miniexammar}
上面和下面的分数非常接近，以至于它们相互接触！ 这不是 Delilah 想要的，她很惊讶 \LaTeX{} 没有检测到这一点并做一些事情。 幸运的是，在 amsmath 环境中，允许将可选参数传递给 \verb=\\= 以定义行之间的实际垂直空间。 对于分数，2ex 是一个不错的选择。 此外，分数处于inline模式。 \verb=\dfrac= 命令给出displayed模式的分数。
\begin{miniexammar}{.3\textandmarginlen}{
\[
\begin{bmatrix}
1&\dfrac{1}{2}&\dfrac{1}{3}\\[2ex]
1&\dfrac{1}{4}&\dfrac{1}{5}
\end{bmatrix}
\]
}
\begin{lstlisting}
\[
\begin{bmatrix}
1&\dfrac{1}{2}&\dfrac{1}{3}\\[2ex]
1&\dfrac{1}{4}&\dfrac{1}{5}
\end{bmatrix}
\]
\end{lstlisting}
\end{miniexammar}

有时矩阵太大而无法完全显示。 在这些时候，使用省略号很重要。 当 Delilah 写出矩阵的逆时，她使用省略号。
\begin{miniexammar}{.5\textandmarginlen}{
\[
A^{-1} = \frac{1}{\det A}
\begin{bmatrix}
C_{11} & C_{21} & \cdots & C_{n1} \\
C_{12} & C_{22} & \cdots & C_{n2} \\
\vdots & \vdots & \ddots & \vdots \\
C_{n2} & C_{n2} & \cdots & C_{nn} \\
\end{bmatrix}
\]
}
\begin{lstlisting}
\[
A^{-1} = \frac{1}{\det A}
\begin{bmatrix}
C_{11} & C_{21} & \cdots & C_{n1} \\
C_{12} & C_{22} & \cdots & C_{n2} \\
\vdots & C\vdots & \ddots & \vdots \\
C_{n2} & C_{n2} & \cdots & C_{nn} \\
\end{bmatrix}
\]
\end{lstlisting}
\end{miniexammar}

\subsection{文字和运算符}
要将文本放入数学环境中，Delilah 使用 amsmath 提供的 \verb=\text= 命令。
\begin{miniexammar}{.57\textandmarginlen}{
\begin{definition}[Null Space]
The null space of an $m \times n$ matrix $A$, written as $\Nul A$, is the set of all solutions
of the homogeneous equation $A\vec{x} = \vec{0}$. In set notation,
\[
\Nul A = \{\vec{x}:\vec{x} \text{ is in } \mathbb{R}^n \text{ and } A\vec{x} = \vec{0} \}
\]
\end{definition}
}
\begin{lstlisting}
\begin{definition}[Null Space]
The null space of an $m \times n$ matrix $A$, written as $\Nul A$, is the set of all solutions
of the homogeneous equation $A\vec{x} = \vec{0}$. In set notation,
\[
\Nul A = \{\vec{x}:\vec{x} \text{ is in } \mathbb{R}^n \text{ and } A\vec{x} = \vec{0} \}
\]
\end{definition}
\end{lstlisting}
\end{miniexammar}

\verb=\Nul=, \verb=\sin=, ... 等命令是数学运算符。 \LaTeX{} 中的部分预定义数学运算符显示在表 \ref{tab:predefmathop} 中。
\begin{table}[hbpt]·
\begin{center}
\small
\begin{tabular}{cl|cl|cl}
Result & Command & Result & Command & Result & Command \\
\hline
arccos & \verb=\arccos= & arcsin & \verb=\arcsin= & arctan & \verb=\arctan= \\
cos & \verb=\cos= & sin & \verb=\sin= & tan & \verb=\tan= \\
cot & \verb=\cot= & sec & \verb=\sec= & csc & \verb=\csc= \\
cosh & \verb=\cosh= & sinh & \verb=\sinh= & tanh & \verb=\tanh= \\
lim & \verb=\lim= & lim inf & \verb=\liminf= & lim sup & \verb=\limsup= \\
ln & \verb=\ln= & log & \verb=\log= & lg & \verb=\lg= \\
max & \verb=\max= & min & \verb=\min= & sup & \verb=\sup= \\
inf & \verb=\inf= &  &  &  &  \\
ker & \verb=\ker= & det & \verb=\det= & exp & \verb=\exp= 
\end{tabular}
\end{center}
\caption{一些预定义的数学运算符}
\label{tab:predefmathop}
\end{table}

实际上，运算符\verb=\Nul= 和\verb=\Span= 是由xjtlumath 定义的，如西浦一年级线性代数教科书中的形式。 此外，xjtlumath 更改了 \LaTeX{} 中的默认 \verb=\vec= 命令，以便矢量以粗体形式出现，而不是在其上方带有箭头。

一些运算符，如 \verb=\lim=，旨在支持对它取极限。 也就是说，在displayed模式下，当试图使用 \verb=_= 给这样的运算符一下标时，下标将出现在运算符的底部。
\begin{parexammar}{.4\textandmarginlen}{
\[
\lim_{x\to 0} f(x)
\]
}
\begin{lstlisting}
\[
\lim_{x\to 0} f(x)
\]
\end{lstlisting}
\end{parexammar}

Delilah 能够通过使用 \verb=\limits= 和 \verb=\nolimits= 来明确控制极限样式。 请注意，这两个命令只能在支持采取极限的操作后使用。
\begin{parexammar}{.4\textandmarginlen}{
$\lim\limits_{x \to 0}f(x)$
\[
\lim\nolimits_{x\to 0} f(x)
\]
}
\begin{lstlisting}
$\lim\limits_{x \to 0}f(x)$
\[
\lim\nolimits_{x\to 0} f(x)
\]
\end{lstlisting}
\end{parexammar}

\subsection{符号}
数学环境中的标准 \LaTeX{} 字体整洁干净。 然而在一些特殊场合，Delilah 想改变一些符号的字体。 例如，为了表示一些常规集合，她使用黑板字体。
\begin{parexammar}{.4\textandmarginlen}{
\[
\mathbb{R}\quad \mathbb{N}\quad \mathbb{Q}\quad \mathbb{Z}
\]
}
\begin{lstlisting}
\[
\mathbb{R}\quad \mathbb{N}\quad \mathbb{Q}\quad \mathbb{Z}
\]
\end{lstlisting}
\end{parexammar}

每次都写 \verb=\mathbb= 有点烦人。 为此，xjtlumath 为它们定义了简写。
\begin{parexammar}{.45\textandmarginlen}{
\[
\setr \quad \setq \quad \setz \quad \setn \quad \setnp
\]
}
\begin{lstlisting}
\[
\setr \quad \setq \quad \setz \quad \setn \quad \setnp
\]
\end{lstlisting}
\end{parexammar}

其他字体控制方法如我们在 \ref{subsec:fonts} 小节中讨论过的。 例如，\verb=\mathrm= 给出\textrm{Roman} 家族的字体，而\verb=\mathbf= 给出\textbf{bold} 系列的字体。
\begin{parexammar}{.45\textandmarginlen}{
\[
\mathrm{Like normal text} \quad \mathbf{bold}
\]
}
\begin{lstlisting}
\[
\mathrm{Like normal text} \quad \mathbf{bold}
\]
\end{lstlisting}
\end{parexammar}

\section{Tutorial 3: Yue Handling Floats}
Yue is preparing her material for the monthly of math club. She wants to make her material interesting and easy to understand, so she utilizes many figures and tables.

Figures, tables and many other things that occupy an (often large) area of random width and height are treated as \emph{floats} in \LaTeX{}. Floats are common in today's documents, yet they cause great troubles in typesetting. In this section, we will work along with Yue to see how floats are handled in \LaTeX{}.

\subsection{Inserting Images}
To insert images in \LaTeX{}, package graphicx (loaded by the template files) is a good option. It provides the \verb=\includegraphics= that accepts the input image filename and some optional specifiers.
\begin{miniexammar}{.4\textandmarginlen}{
\includegraphics[width=\textwidth]{assets/examplelogo.jpg}
}
\begin{lstlisting}
\includegraphics[width=\textwidth]{assets/examplelogo.jpg}
\end{lstlisting}
\end{miniexammar}

Yue soon finds out that simply using this command is not a good option, because if there is not enough vertical space for the image, it will be placed on the next page, leaving a large vacant area, which is quite ugly. Also, she is unable to provide the image with a caption or to reference it.

So, Yue wraps the image with the figure environment that makes the image a \emph{figure}.
\begin{miniexammar}{.45\textandmarginlen}{
% figure with a caption cannot be placed inside a minipage, so we fake it here. 
\includegraphics[width=\textwidth]{assets/examplelogo.jpg}
\begin{center}
\hypertarget{fakedcaption}{Figure 1: Example Logo}
\end{center}
Figure \hyperlink{fakedcaption}{1} shows the figure Yue uses.
}
\begin{lstlisting}
\begin{figure}
\includegraphics[width=\textwidth]{assets/examplelogo.jpg}
\caption{Example Logo}
\label{fig:example}
\end{figure}
Figure \ref{fig:example} shows the figure Yue uses.
\end{lstlisting}
\end{miniexammar}

\LaTeX{} automatically gives a number to it so that Yue is able to reference it. Note that due to implementation reasons, \verb=\label= should only be placed immediately after a \verb=\caption= lest the reference be wrong.

\subsection{Tables}
Even if using figures does require an extra environment, it is still quite simple and Yue soon becomes familiar with it. Dealing with tables, however, is more complex in \LaTeX{}.

To generate a table-like content in \LaTeX, Yue has to use special environments. tabular and array are two specific examples of them. In fact, the two environments are alike in most of the aspects, with one major difference being that array is often used in math mode.

The syntax of array and tabular resembles the one of matrix environments used by Delilah, though here Yue has to explicitly specify the column behavior.
\begin{miniexammar}{.4\textandmarginlen}{
\begin{tabular}{|c|c|}
\hline 
Entry 1 & Entry 2\\
\hline
a & b\\
\hline
\end{tabular}
}
\begin{lstlisting}
\begin{tabular}{|c|c|}
\hline 
Entry 1 & Entry 2\\
\hline
a & b\\
\hline
\end{tabular}
\end{lstlisting}
\end{miniexammar}

Yue doesn't quite understand what \verb=|c|c|= means, so she searches on the Internet for this. It tells her that this argument passed to tabular specifies each column. The letter c tells tabular that the contents in this column should be centered. And two other alignment specifier l and r are available for ``left'' and ``right'', respectively. The vertical bar indicates that at this place a vertical line should be inserted (c.f. The \verb=\hline= command that instructs a horizontal line to be inserted at the top of the current row)

The width of a column is determined by the width of the contents when either c,l, and r is given. Yue is able to control the width of a column by using another specifier ``p'', in which the content is left-aligned.
\begin{miniexammar}{.4\textandmarginlen}{
\begin{tabular}{|c|p{2cm}|}
\hline 
Entry 1 & Entry 2\\
\hline
a & b\\
\hline
\end{tabular}
}
\begin{lstlisting}
\begin{tabular}{|c|p{2cm}|}
\hline 
Entry 1 & Entry 2\\
\hline
a & b\\
\hline
\end{tabular}
\end{lstlisting}
\end{miniexammar}

When there is many columns with the same specifier, Yue can use this syntax \verb=*{num}{spe}= to repeat the specifiers, where num is the number of repetitions, and spe is the specifier.
\begin{miniexammar}{.4\textandmarginlen}{
\begin{tabular}{|*{7}{c|}}
\hline 
a&a&a&a&a&a&a\\
\hline
a&a&a&a&a&a&a\\
\hline
\end{tabular}
}
\begin{lstlisting}
\begin{tabular}{|*{7}{c|}}
\hline 
a&a&a&a&a&a&a\\
\hline
a&a&a&a&a&a&a\\
\hline
\end{tabular}
\end{lstlisting}
\end{miniexammar}

Yue doesn't like line-separated tables because she considers them not tidy. She wants to use space to separate contents. She can use \verb=@{\hspace{}}= between column specifiers to specify the inter-column space and use \verb=\vspace{}= immediately before \verb=\\= to add extra space before the next row.
\begin{miniexammar}{.4\textandmarginlen}{
\begin{tabular}{c@{\hspace{1cm}}cc}
a & b & b\\
a & b & b\vspace{.5cm}\\
c & d & d\vspace{.5cm}\\
\end{tabular}
}
\begin{lstlisting}
\begin{tabular}{c@{\hspace{1cm}}cc}
a & b & b\\
a & b & b\vspace{.5cm}\\
c & d & d\vspace{.5cm}\\
\end{tabular} 
\end{lstlisting}
\end{miniexammar}
In fact, inside \verb=@{}= Yue can use not only spaces, but any other contents.
\begin{miniexammar}{.45\textandmarginlen}{
\begin{tabular}{c@{ <POLICE, stay away> }c}
crime scene & people \\
crime scene & people \\
\end{tabular}
}
\begin{lstlisting}
\begin{tabular}{c@{ <POLICE, stay away> }c}
crime scene & people \\
crime scene & people \\
\end{tabular}
\end{lstlisting}
\end{miniexammar}

Like the environment figure, the environment table is designed to accept a table-like content.
\begin{miniexammar}{.4\textandmarginlen}{
\begin{tabular}{cc}
a & b \\
c & d \\
\end{tabular}
\begin{center}
\hypertarget{fakedcaptiontab}{Table 1: Example Table}
\end{center}
Table \hyperlink{fakedcaptiontab}{1} shows a table.
}
\begin{lstlisting}
\begin{table}
\begin{tabular}{cc}
a & b \\
c & d \\
\end{tabular}
\caption{Example Table}
\label{tab:example}
\end{table}
Table \ref{tab:example} shows a table.
\end{lstlisting}
\end{miniexammar}

Now that Yue understands how to manipulate table-like contents in \LaTeX{}, she soon resumes to her work. At some point, she has to create a super big table that contains at least 20 columns, and she finds that it sticks out of the page. When she tries to use \verb=\small= to reduce the size of the texts, she finds that she has to copy it for every entry of the table, which means hundreds times of repetition. To control the overall style, Yue needs to give the controlling commands between the beginning of table and the beginning of tabular.

\begin{miniexammar}{.5\textandmarginlen}{
{
\scriptsize
\begin{tabular}{|*{14}{c|}}
\hline 
a&a&a&a&a&a&a&a&a&a&a&a&a&a\\
\hline
a&a&a&a&a&a&a&a&a&a&a&a&a&a\\
\hline
\end{tabular}
}
}
\begin{lstlisting}
\begin{table}
\small
\begin{tabular}{|*{14}{c|}}
\hline 
a&a&a&a&a&a&a&a&a&a&a&a&a&a\\
\hline
a&a&a&a&a&a&a&a&a&a&a&a&a&a\\
\hline
\end{tabular}
\end{table}
\end{lstlisting}
\end{miniexammar}

\subsection{Placement of Floats}
Yue used to use Microsoft Word, which places floats wherever the user wants to place it at. Everything had been fine since she had turned to \LaTeX for some time, but now Yue has a problem. A figure ``disappeared'' from the output. After checking her code and the output again and again, she accidentally finds that the figure appears in the next page. This really confuses her. Due to the internal algorithm of \TeX, it is technically impossible to arrange every float at where the user wants to place them. According to \textit{the \LaTeX{} Companion},
\begin{quotation}
``Floats are often problematic in the present version of \LaTeX, because the system was developed at a time when documents contained considerably less graphical material than they do today.''
\end{quotation}

Yet \LaTeX{} does give some option to allow the Yue to control the placement of a float to some extend. For a figure or table environment, Yue is able to pass to it an optional argument that specify the desired placement. There are five placement specifiers and they can be combined together in any order.
\begin{description}
\item[!] Ignores some \LaTeX{} limitation\footnote{\LaTeX{} has some limitations when it tries to place floats. For example, if the height of a float is larger than a degree of page height, then it cannot be placed at the bottom of a page.} when trying to place the float.
\item[h] Tries to place the float exactly at where the the environment is issued. If the attempt fails and no other specifier other that ! is given, the specifier will be changed to t.
\item[t] Tries to place the float at the top of the page.
\item[b] Tries to place the float at the bottom of the page.
\item[p] Tries to place the float at a float page (a page that is generated by \LaTeX{} to place floats)
\end{description}
\LaTeX{} tries to place a float according to the specifiers in the order of the above list from the top to the bottom. Generally all floats of a document can be handled properly. But should a float proves impossible to handle, the author should adjust (probably reduce) its width and height.

\subsection{Table of Floats}
As mentioned at the beginning, in Yue's material, there are many floats. She wonders if there is a way to provide a quick reference to them.

Like \verb=\tableofcontents=, \LaTeX{} provides the following two commands to print a list of all figures and a list of all tables used in a document, respectively.
\begin{verbatim}
\listoffigures and \listoftables
\end{verbatim}
The name of a float that appears on the list is defined by the caption of the float. If the captions appears to be too long, Yue can pass to caption an optional argument that will instead be shown on the list. Also, don't forget to compile the file at least twice for the lists to show properly.

\subsection{A Suggestion About Images}
Yue is suggested to use vector graph for images, as vector graph is lossless when the image is scaled along with the output file.

There are several packages that enables drawing images directly in \LaTeX{}, yet they all require great efforts. It is suggested to use modern tools to generate appropriate images (e.g. Mathematica is able to export the math plots drawn.).

\section{Tutorial 4: The Story of ZiYou and Abigail} % ZiYou for 子由
ZiYou and Abigail are the team leaders of the material about Calculus for the final exam. As team leaders, they have to deal with more problems than their colleagues. In this section, we will know about how ZiYou and Abigail manage to solve these problems and how their affection of each other grows.

\subsection{Managing Notes In \LaTeX}
In a review to the work of one of the team members, ZiYou finds several places that may not be clear enough for the readers. He decides to add some description to these unclear texts. These descriptions should not defer the reading of the main text, so ZiYou adjudicates on making them \emph{notes}. Two general ways for adding notes in \LaTeX{} are using \emph{footnotes} and using \emph{marginpars}.

A footnote in \LaTeX{} provides an annotation for a piece of text in the footer of the current page, and generates a number of the note which will appear as the superscript of the text being annotated. To use footnotes, ZiYou uses the command \verb=\footnote=.
\begin{miniexammar}{.55\textandmarginlen}{
This is something unclear\footnote{This means that ...}. And some other texts are here.
}
\begin{lstlisting}
This is something unclear\footnote{This means that ...}. And some other texts are here.
\end{lstlisting}
\end{miniexammar}

Different from ZiYou, Abigail prefers to use marginpars for annotation. A marginpar appears in the margin of the current page, but does not possess a number like a footnote does.
\begin{miniexammar}{.6\textandmarginlen}{
% We have to fake a margin par here.
\parbox{.68\textwidth}{This is something unclear. And some other texts are here.} \hspace{.03\textwidth}
\parbox{.27\textwidth}{That is, we have to ...}
}
\begin{lstlisting}
This is something unclear\marginpar{That is, we have to ...}. And some other texts are here.
\end{lstlisting}
\end{miniexammar}

A marginpar appears in the margin, and is at the same height as the text where the marginpar is given. When Abigail sees that ZiYou uses footnotes rather than marginpars, she asks him to change them because she thinks that marginpars are better. Certainly ZiYou doesn't agree with her, but he confers with what she claims, that the noting styles should agree with each other in a document.

To give a resolution about what noting style is to be used, ZiYou suggests that they play Tic Tac Toe, and Abigail thinks this is a good idea. After a few minutes, Abigail narrowly wins the game. ZiYou jokes that maybe she should let him win the next time, whereupon Abigail smilingly replies, ``That remains to be seen''. Nevertheless, she has a dim feeling that ZiYou deliberately lets her be the winner, but can't prove it from ZiYou's regretful expression.

\subsection{Merging Works Of The Team}
ZiYou and Abigail only need to collect the chapter.tex from the team members. They rename the files according to each person in a way that they can easily identify who is responsible for each file. After that, they input each file into the chapter.tex of an empty template. The final output can then be generated from the encapsulation.tex of the template.

Directly copying the contents is not a good option for inputting the files. ZiYou is about to search this on the Internet when Abigail discovers in the encapsulation.tex that the chapter.tex is directed into this file by the command \verb=\input=. The following code shows what is written in encapsulation.tex.
\begin{lstlisting}
\input{chapter.tex}
\end{lstlisting}
The argument passed to this command is the relative path of the target file. Abigail doesn't know what the term relative path means, so ZiYou, who has certain knowledge in computer science, explains to her that, the relative directory is the path of a file relative to the file in which the path is used. In this example, the file in which the relative directory is used is encapsulation.tex, and as a consequence of the two files being in the same folder, the relative directory of the chapter.tex is simply its name. If ZiYou and Abigail decide to put the files of the team members into a folder named Files for organization, then the directory should contain the folder's name plus a \verb=/= or \verb=\=, depending on the file system, at the beginning. Since that \verb=\= is a keyword in \LaTeX, they should use \verb=/=.

Not until they have completed the inputting had ZiYou had a glimpse at the search result on the Internet, where he finds another command, \verb=\include=. After seeing this webpage in detail, he then tells Abigail that \verb=\include= is a better choice here, for it somehow improves the compilation speed. Abigail has no clue about what compilation is and isn't interested in such technical stuffs, but she trusts ZiYou. She also kind of likes it when ZiYou, patiently and tenderly, explains what she doesn't understand to her. So she pretends to be curious and asks ZiYou to explain compilation to her.

So the final form of their chapters.tex is like this:
\begin{lstlisting}
\include{Files/The first file}

\include{Files/The second file}

\include{Files/The third file}
...
\end{lstlisting}
Note that the file extension (.tex) is not allowed to be used in \verb=\include=, while it can be used in \verb=\input=. Also, they have to make sure that the team members have not taken advantage of \verb=\include=, as it cannot be used in a file that is included by another. Fortunately, they can ascertain it since none of them knows about the command.

\subsection{Background, Headers, and Footers}
ZiYou and Abigail notice that the template for materials automatically adds the background, and the header and footer for each page. The encapsulation.tex loads the package background for background, and the package fancyhdr for headers and footers.

The headers and footers are set in the encapsulation.tex by the following code: (A line begins with a \verb=%= is \emph{commented}, so it will have no effect on the output .pdf file.)
\begin{lstlisting}
% Define the header and footer for pages.

% Place the number of the current page.
\fancyhead[LEH,ROH]{\bfseries\thepage}

% Beautify the display of chapter and section marks.
\renewcommand{\chaptermark}[1]{%
\markboth{#1}{}}
\renewcommand{\sectionmark}[1]{%
\markright{\thesection\ #1}}

\fancyhead[LOH]{\bfseries\rightmark}
\fancyhead[REH]{\bfseries\leftmark}

% Add copyright in the footer
\fancyfoot[COF,CEF]{\bfseries\copyright{} The XJTLU Math Club -- All rights reserved}
\end{lstlisting}
The header is controlled by the command \verb=\fancyhead= and the footer is controlled by the command \verb=\fancyfoot=. By inspecting the optional arguments, ZiYou makes a guess that L and R represent left and right, E and O represent even and odd (page number), and H and F represent header and footer, respectively. A look into the documentation of fancyhdr confirms this. He asks Abigail if she likes the page style. Abigail thinks that the author has a good taste, so they decide not to change this.

The background is set to the logo of math club. In fact, adding this graph somehow makes the document ugly, and even I didn't understand why this must be added to all materials. The head of the department at that time told me that this is a defense for those who use the materials in a prohibited way.

I really hope that this environment can improve to a state that even the materials are distributed by the most free licenses, there will be no one to steal our intelligent property.

\subsection{Index And Bibliography}
ZiYou and Abigail would like to listen to the readers' opinions about the previous materials so that they can refine the coming one according to them. A few number of the readers pointed out that it took them a lot of efforts to find specific terms in the materials and they would appreciate it if a list of important terms is added.

Abigail recalls that once when she tried to find something in a calculus textbook's appendix, she flipped the pages too much and turned to a section called Index, where terms are listed according to their pages. So ZiYou makes a search on the Internet and found that Indexing in \LaTeX{} can be easily done by a few commands.

First, for any important terms that they want to list in the Index page, they mark it with the command \verb=\index=. To print the Index page, the command \verb=\printindex= needs to be called, and before the document environment the command \verb=\makeindex= should be called. In the encapsulation.tex, simply uncomment the relating lines of code to do that. Finally, for the Index page to be printed, they have to, first, run \LaTeX{} once on encapsulation.tex, then, run MakeIndex once on the file, and finally, run \LaTeX{} twice on the file. The following example show the result of using Index.
\begin{miniexammar}{.35\textandmarginlen}{
\section*{Index}
limit, vi, 3\\
derivative, 3-5
}
\begin{lstlisting}
% At page vi
\index{limit}
% At page 3
\index{limit}
\index{derivative}
% At page 4
\index{derivative}
% At page 5
\index{derivative}
\end{lstlisting}
\end{miniexammar}

ZiYou makes definite integral and indefinite integral two separate index entries, while Abigail argues that they should be under the same term integral. To use subterms, the following syntax should be applied.
\begin{miniexammar}{.4\textandmarginlen}{
\section*{Index}
integral, 5
\par \hspace{2em} definite integral, 11
\par \hspace{2em} indefinite integral, 7
}
\begin{lstlisting}
% At page 5
\index{integral}
% At page 7
\index{integral!definite integral}
% At page 11
\index{integral!indefinite integral}
\end{lstlisting}
\end{miniexammar}

Abigail remembers that for the sake of academic integrity, they should add reference for each work of other people. B\textsc{ib}\TeX{} is a good tool handling references (bibliographies). To use this tool, they have to prepare a B\textsc{ib}\TeX{} database. This is a file with the extension .bib, in which entries are contained. Filling in the entries is a tedious work, but fortunately most academic websites provides the facility to allow an user to directly download a .bib file for the source he wants to use. The following code shows a typical database entry:
\begin{lstlisting}
@article{may1979alpha,
	title="Alpha-particle-induced soft errors in dynamic memories",
	author="T.C. {May} and M.H. {Woods}",
	journal="IEEE Transactions on Electron Devices",
...
}
\end{lstlisting}
ZiYou and Abigail need not to worry about the details in the entry. The only thing they need to remember is the label of the entry, the one immediately after the \verb={=, because it is to be used in the \verb=\cite= command to produce a reference to that source. Each source that is referenced in the document appears in the Bibliography page, which is controlled by the two following commands:
\begin{lstlisting}
\bibliography{file-list}
\bibliographystyle{style}
\end{lstlisting}
, where file-list is a list of database files, and style is the bibliography style according to which the bibliography is to be printed. \url{https://www.overleaf.com/learn/latex/Bibtex_bibliography_styles} shows all predefined B\textsc{ib}\TeX{} styles.

\subsection{Twosided printing}
Now ZiYou and Abigail have finished combining the separate documents and are planning a break after the work, but there is something more waiting for them.

By default, when the documentclass is book, \LaTeX{} uses twosided printing. At first ZiYou and Abigail didn't notice this (they didn't know) until Abigail saw that the pages of the output are not aligned. That is, at one page the content sticks out to the left and at the next page the content sticks out to the right. Abigail doesn't want to trouble ZiYou again, so she makes a search herself and finds that this is what is called twoside, which is used particularly for books.

Abigail tries to comprehend this by opening a book and inspecting its layout. She notices that a portion of the inner part\footnote{When one opens a book, the first page appearing on the right is numbered 1, and all odd pages appears on the right side have odd numbers as a consequence.} of the pages are stuck together to keep them form a book. She guesses that twoside leaves extra spaces for each page at the inner part. The result, however, contradicts with her instinct. On twoside mode, \LaTeX{} shrinks the inner part of each page. This is because \LaTeX{} wants to give more space for the margins that reside on the outer parts.

ZiYou learns this from Abigail so they start another review for it. They soon find a problem. For table of contents and some other pages like preface, the template uses roman numbering. For the main text, the page numbering is changed to arabic. The layout of a page, however, is not decided by if it will appear on the left or right. Rather, it is entirely decided by the number of the page. If, say, the roman number ends at iii, then when the arabic number starts at 1, the corresponding page will appear on the left, which is a disaster. To solve this problem, ZiYou finds the command \verb=\cleardoublepage= that is designed for twoside mode. It ensures that the content after the command will appear at a page on the right by (possibly) adding a new page. The new page would be empty, of course. But Abigail thinks that a reader may misinterpret this as a sign of a printing failure. To explicitly identify this behavior, in xjtlumath (in precise, the subpackage xjtlumaterial loaded by the template), the command is renewed to add the sentence ``This page is intentionally left blank'' at each empty page this command adds.

\subsection{An End Of Their Story}
The very story of ZiYou and Abigail about their material writing comes to an end, yet as said in one old saying, ``An end is also a beginning,'' their other stories just begin. Abigail admires ZiYou's broad knowledge in computer science and enjoys it when ZiYou teaches her about it. ZiYou also is inspired by Abigail's energy, particularly in the working. He feels lucky to find such a girl as interesting and beautiful as Abigail, whose humor has made ZiYou laughing several times.

So the two ordered seats in a romantic restaurant to celebrate their collaboration, while I, sitting alone in my dormitory, work hard to complete this documentation. Nevertheless, I hope you, the members of the math club in the future, can find not only knowledge and experience, but also love as ZiYou and Abigail do, except not in an imaginary story.

\part{模板}
这部分描述了所有给出的模板。 它们位于文件夹 \verb|Templates| 中。
\section{资料专用模板}
在本节中，所有模板都是为写作资料制作的。 它们位于文件夹 \verb|For Materials| 中。
\subsection{英语书籍}
文件夹 \verb=Book_en-us= 提供了用于编写将作为书籍出版的英文材料的模板。 长材料应以本模板的形式准备。 它加载所有 xjtlumath 包，因此该包提供的所有功能都可用。

在环境document的开头，页码改为罗马，页面样式设置为plain，并给出了目录。 您可以在此之前添加标题页。 此外，其他前言，如致谢和dedication，也应包括在此处。

从那里开始正文，其中页码是阿拉伯语，样式是fancy，这是由之前的几个fancyhdr命令定义的。

在文档环境的末尾，您可以取消注释一些命令以使用索引和参考书目（索引开头需要一个命令）。

\subsection{中文书籍}
文件夹\verb=Book_zh-cn= 提供了用于编写将作为书籍出版的中文材料的模板。 除了文档类是 ctex 包提供的 ctexbook 中文样式外，与上一个相同。

\subsection{英语文章}
文件夹 \verb=Article_en-us= 提供了用英语编写材料的模板，该材料将作为文章发表。 大小适中的材料应以本模板的形式准备。 它加载所有 xjtlumath 包，因此该包提供的所有功能都可用。

它类似于英文书籍的模板，只是将twoside的一些选项调整为适合oneside。

\subsection{中文文章}
文件夹\verb=Article_zh-cn= 提供了用于编写将作为文章发表的中文材料的模板。 除了文档类是 ctex 提供的 ctexart 中式样式外，与上一个相同。

\section{一般使用}
在本节中，所有模板都是为一般用途制作的。 它们位于文件夹 \verb|General Use| 中。

目前它们与材料的模板相同，只是它们不加载包 xjtlumaterial 并且没有一些样式。


\part{内部文档}
这部分是 xjtlumath 的内部文档。 它由工具doc自动生成。 这部分用于包的维护，不建议包的一般用户阅读，除非他们想了解包的实现细节。很抱歉内部文档注释不包含中文版。
\DocInput{xjtlumath.dtx} 

\clearpage
\DocInput{mydoc.dtx}
\end{document}